\documentclass[11pt]{article}
\usepackage[margin=1in]{geometry}
\usepackage{hyperref}
\usepackage[T1]{fontenc}
\usepackage[utf8]{inputenc}
\title{Discovery Branch Synthesis}
\author{Abducer}
\date{}
\begin{document}
\maketitle

\noindent
\textbf{KB ref:} p-vs-np\\
\textbf{Completed branches:} 0\\
\textbf{Executed branch passes:} 25\\
\textbf{Proof queue pending:} 71\\
\textbf{Proof queue completed:} 38\\
\textbf{Ready for prove:} 2\\
\textbf{Ready for discovery:} 23

\section*{Judgment}
\begin{itemize}
\item All routed discovery branches were executed once. The frontier is now better characterized than expanded: the dominant pressure remains on MCSP carrier realization, formula-size transport, and theorem-activation alignment.
\item Next move: Promote only the theorem-shaped branch families into the next prove-or-disprove agenda, and keep the pattern-level branches as discovery guidance rather than treating them as solved proof tasks.
\end{itemize}

\section*{Top Breakthrough Families}
\begin{itemize}
\item No repeated breakthrough families recorded.
\end{itemize}

\section*{Top Obstruction Classes}
\begin{itemize}
\item No repeated obstruction classes recorded.
\end{itemize}

\section*{Most Reused InformalMathLib References}
\begin{itemize}
\item Hardness Magnification Near State-of-the-Art Lower Bounds: 15
\item Equivalence and coherent transport theorem for MCSP witness family: 11
\item Hardness Magnification for Natural Problems: 11
\item A Simple and General Hardness Magnification Theorem: 7
\item Polynomially Bounded Proof Systems Characterize NP Equals coNP: 6
\item MCSP witness family: major consequences such as NP not in P/poly: 5
\item Hitting-Set Generators Yield Black-Box PIT: 5
\item Commutative factorization theorem for MCSP obstruction family: 4
\end{itemize}

\section*{Ready For Prove}
\begin{itemize}
\item disclayer.p-vs-np.family.barrier-aware-transfer: There exists an admissible 'unspecified witness family' such that 'There exist many problems L in combinatorial optimisation/graph theory such that L is NP-complete under polynomial-time many-one reductions.' can be organized into a route preserving 'consequence schema NP-complete-classification' and thereby supports 'Under"local"polynomial-time reductions (where each output bit is computable with very small time), MCSP is not NP-complete. More broadly, the paper develops evidence and conditional consequences suggesting NP-hardness of MCSP under various reductions would entail unexpected class separations.'. (priority 2.309; linked queue 3)
\item disclayer.p-vs-np.family.explicitness: There exists an admissible 'PIT certificate family' such that 'A set Hsubseteq mathbb F \textasciicircum{}n is a hitting set for a class mathcal C of n -variate circuits if for every nonzero polynomial p computed by some circuit in mathcal C , there exists mathbf a in H with p(mathbf a )neq 0.' can be organized into a route preserving 'PIT semantic object' and thereby supports 'If mathrm PIT in Pclass (or even suitable infinitely-often derandomisations in nondeterministic subexponential time, in standard variants), then at least one of the following holds: NEXP not in P/poly quad or quad Permanent does not have polynomial-size arithmetic circuits.'. (priority 2.107; linked queue 3)
\end{itemize}

\section*{Ready For Discovery}
\begin{itemize}
\item disclayer.p-vs-np.pattern.invariant-discovery: There exists an admissible 'MCSP witness family' such that 'A hardness magnification theorem has the form:"A barely-nontrivial lower bound for an explicit problem Q in a weak model implies a major lower bound/separation in a strong model."Typically, Q is a meta-problem (e.g. a Gap-MCSP variant) and"major"means consequences such as NP not in P/poly.' can be organized into a route preserving 'MCSP complexity profile in nonuniform-circuits' and thereby supports 'In refined regimes, lower bounds for Gap-MCSP slightly beyond known frontiers (e.g. N\textasciicircum{} 1+varepsilon -type circuit lower bounds for suitable promise gaps and input length N=2\textasciicircum{}n ) imply major consequences such as NP not in P/poly.'. (priority 6.868; linked queue 10)
\item disclayer.p-vs-np.pattern.invariant-discovery: There exists an admissible MCSP witness family, realized as a family of instances equipped with a monotone witness measure, such that whenever A barely-nontrivial lower bound for an explicit problem Q in a weak model holds for MCSP in nonuniform-circuits, the object organizes the source-side mathematical content while preserving the invariant MCSP complexity profile in nonuniform-circuits and yields major consequences such as NP not in P/poly. (priority 6.868; linked queue 10)
\item disclayer.p-vs-np.pattern.regime-localization: There exists an admissible 'MCSP witness family' such that 'A hardness magnification theorem has the form:"A barely-nontrivial lower bound for an explicit problem Q in a weak model implies a major lower bound/separation in a strong model."Typically, Q is a meta-problem (e.g. a Gap-MCSP variant) and"major"means consequences such as NP not in P/poly.' can be organized into a route preserving 'MCSP complexity profile in nonuniform-circuits' and thereby supports 'In refined regimes, lower bounds for Gap-MCSP slightly beyond known frontiers (e.g. N\textasciicircum{} 1+varepsilon -type circuit lower bounds for suitable promise gaps and input length N=2\textasciicircum{}n ) imply major consequences such as NP not in P/poly.'. (priority 6.758; linked queue 10)
\item disclayer.p-vs-np.pattern.regime-localization: There exists an admissible MCSP witness family, realized as a family of instances equipped with a monotone witness measure, such that whenever A barely-nontrivial lower bound for an explicit problem Q in a weak model holds for MCSP in nonuniform-circuits, the object organizes the source-side mathematical content while preserving the invariant MCSP complexity profile in nonuniform-circuits and yields major consequences such as NP not in P/poly. (priority 6.758; linked queue 10)
\item disclayer.p-vs-np.pattern.hybridization-of-partial-bridges: There exists an admissible 'MCSP witness family' such that 'A hardness magnification theorem has the form:"A barely-nontrivial lower bound for an explicit problem Q in a weak model implies a major lower bound/separation in a strong model."Typically, Q is a meta-problem (e.g. a Gap-MCSP variant) and"major"means consequences such as NP not in P/poly.' can be organized into a route preserving 'MCSP complexity profile in nonuniform-circuits' and thereby supports 'In refined regimes, lower bounds for Gap-MCSP slightly beyond known frontiers (e.g. N\textasciicircum{} 1+varepsilon -type circuit lower bounds for suitable promise gaps and input length N=2\textasciicircum{}n ) imply major consequences such as NP not in P/poly.'. (priority 6.728; linked queue 10)
\item disclayer.p-vs-np.pattern.hybridization-of-partial-bridges: There exists an admissible MCSP witness family, realized as a family of instances equipped with a monotone witness measure, such that whenever A barely-nontrivial lower bound for an explicit problem Q in a weak model holds for MCSP in nonuniform-circuits, the object organizes the source-side mathematical content while preserving the invariant MCSP complexity profile in nonuniform-circuits and yields major consequences such as NP not in P/poly. (priority 6.728; linked queue 10)
\item disclayer.p-vs-np.pattern.counterexample-led-concept-formation: There exists an admissible MCSP witness family, realized as a family of instances equipped with a monotone witness measure, such that whenever A barely-nontrivial lower bound for an explicit problem Q in a weak model holds for MCSP in nonuniform-circuits, the object organizes the source-side mathematical content while preserving the invariant MCSP complexity profile in nonuniform-circuits and yields major consequences such as NP not in P/poly. (priority 6.658; linked queue 10)
\item disclayer.p-vs-np.pattern.route-decomposition: There exists an admissible MCSP witness family, realized as a family of instances equipped with a monotone witness measure, such that whenever A barely-nontrivial lower bound for an explicit problem Q in a weak model holds for MCSP in nonuniform-circuits, the object organizes the source-side mathematical content while preserving the invariant MCSP complexity profile in nonuniform-circuits and yields major consequences such as NP not in P/poly. (priority 6.578; linked queue 10)
\item disclayer.p-vs-np.pattern.problem-representation-shift: There exists an admissible MCSP witness family, realized as a family of instances equipped with a monotone witness measure, such that whenever A barely-nontrivial lower bound for an explicit problem Q in a weak model holds for MCSP in nonuniform-circuits, the object organizes the source-side mathematical content while preserving the invariant MCSP complexity profile in nonuniform-circuits and yields major consequences such as NP not in P/poly. (priority 6.548; linked queue 10)
\item disclayer.p-vs-np.family.realization: If source theorem S holds in regime R, then one can construct an explicit carrier C with admissibility conditions A such that C preserves invariant I and supports a route-specific forcing lemma to target T. (priority 6.443; linked queue 9)
\item disclayer.p-vs-np.family.abstraction-search: There exists an admissible 'MCSP witness family' such that 'A hardness magnification theorem has the form:"A barely-nontrivial lower bound for an explicit problem Q in a weak model implies a major lower bound/separation in a strong model."Typically, Q is a meta-problem (e.g. a Gap-MCSP variant) and"major"means consequences such as NP not in P/poly.' can be organized into a route preserving 'MCSP complexity profile in nonuniform-circuits' and thereby supports 'In refined regimes, lower bounds for Gap-MCSP slightly beyond known frontiers (e.g. N\textasciicircum{} 1+varepsilon -type circuit lower bounds for suitable promise gaps and input length N=2\textasciicircum{}n ) imply major consequences such as NP not in P/poly.'. (priority 6.443; linked queue 9)
\item disclayer.p-vs-np.family.abstraction-search: State a route-specific theorem that would resolve the blocked interface and make the next proof obligation materially sharper. (priority 6.443; linked queue 9)
\item disclayer.p-vs-np.family.bridge-specificity: Target theorem T is fixed. The real theorem under search is the local activation theorem A that verifies T's hypotheses from current artifacts. (priority 5.669; linked queue 9)
\item disclayer.p-vs-np.family.transport: For every admissible carrier C in the nonuniform-circuit regime, there is a transformed carrier C' in the formula-size regime preserving invariant I and retaining enough quantitative strength to imply target T. (priority 5.610; linked queue 8)
\item disclayer.p-vs-np.family.barrier-aware-transfer: Either construct a local reduction family G satisfying locality bounds L and preserving witness structure W, or prove that any such G would contradict barrier theorem B. (priority 2.309; linked queue 3)
\item disclayer.p-vs-np.family.alignment: Every admissible carrier C from the local route induces an object D in the target theorem's class, and D preserves the quantitative property needed to activate the target theorem. (priority 2.222; linked queue 3)
\item disclayer.p-vs-np.family.explicitness: If semantic object C exists for the relevant circuit class, then one can uniformly and efficiently extract an explicit structure G that yields deterministic PIT for that class. (priority 2.107; linked queue 3)
\item disclayer.p-vs-np.family.constructive-instantiation: From abstract schema S, construct an explicit family of natural instances E\_n satisfying properties P and realizing target theorem T. (priority 1.966; linked queue 3)
\item disclayer.p-vs-np.family.semantic-bridge: Local definitional object D has property P if and only if global characterization condition Q holds, with both directions proved at the right level of abstraction. (priority 1.946; linked queue 3)
\item disclayer.p-vs-np.family.categorical-reframing: There exists a mediating object built from 'MCSP witness family' such that the route from 'A hardness magnification theorem has the form:"A barely-nontrivial lower bound for an explicit problem Q in a weak model implies a major lower bound/separation in a strong model."Typically, Q is a meta-problem (e.g. a Gap-MCSP variant) and"major"means consequences such as NP not in P/poly.' to 'In refined regimes, lower bounds for Gap-MCSP slightly beyond known frontiers (e.g. N\textasciicircum{} 1+varepsilon -type circuit lower bounds for suitable promise gaps and input length N=2\textasciicircum{}n ) imply major consequences such as NP not in P/poly.' factors through typed morphisms preserving 'MCSP complexity profile in nonuniform-circuits' and making the relevant square commute. (priority 0.910; linked queue 12)
\item disclayer.p-vs-np.family.categorical-reframing: There exists a mediating object built from 'MCSP witness family' such that the route from 'A hardness magnification theorem has the form:"A barely-nontrivial lower bound for an explicit problem Q in a weak model implies a major lower bound/separation in a strong model."Typically, Q is a meta-problem (e.g. a Gap-MCSP variant) and"major"means consequences such as NP not in P/poly.' to 'In refined regimes, lower bounds for Gap-MCSP slightly beyond known frontiers (e.g. N\textasciicircum{} 1+varepsilon -type circuit lower bounds for suitable promise gaps and input length N=2\textasciicircum{}n ) imply major consequences such as NP not in P/poly.' factors through typed morphisms preserving 'MCSP complexity profile in nonuniform-circuits' and making the relevant square commute. (priority 0.910; linked queue 12)
\item disclayer.p-vs-np.family.equivalence-coherence-reframing: There exists an equivalence witness identifying 'MCSP witness family' with the target-side object class relevant to 'In refined regimes, lower bounds for Gap-MCSP slightly beyond known frontiers (e.g. N\textasciicircum{} 1+varepsilon -type circuit lower bounds for suitable promise gaps and input length N=2\textasciicircum{}n ) imply major consequences such as NP not in P/poly.', and 'MCSP complexity profile in nonuniform-circuits' transports coherently along that witness. (priority 0.880; linked queue 12)
\item disclayer.p-vs-np.family.equivalence-coherence-reframing: There exists an equivalence witness identifying 'MCSP witness family' with the target-side object class relevant to 'In refined regimes, lower bounds for Gap-MCSP slightly beyond known frontiers (e.g. N\textasciicircum{} 1+varepsilon -type circuit lower bounds for suitable promise gaps and input length N=2\textasciicircum{}n ) imply major consequences such as NP not in P/poly.', and 'MCSP complexity profile in nonuniform-circuits' transports coherently along that witness. (priority 0.880; linked queue 12)
\end{itemize}

\section*{Next Programmes}
\subsection*{disclayer.p-vs-np.pattern.invariant-discovery}
\begin{itemize}
\item Mode: discovery-candidate
\item Target: Invariant Discovery Prompt
\item Routing priority: 6.868
\item Linked queue items: 10
\item Headline claim: There exists an admissible 'MCSP witness family' such that 'A hardness magnification theorem has the form:"A barely-nontrivial lower bound for an explicit problem Q in a weak model implies a major lower bound/separation in a strong model."Typically, Q is a meta-problem (e.g. a Gap-MCSP variant) and"major"means consequences such as NP not in P/poly.' can be organized into a route preserving 'MCSP complexity profile in nonuniform-circuits' and thereby supports 'In refined regimes, lower bounds for Gap-MCSP slightly beyond known frontiers (e.g. N\textasciicircum{} 1+varepsilon -type circuit lower bounds for suitable promise gaps and input length N=2\textasciicircum{}n ) imply major consequences such as NP not in P/poly.'.
\item Rationale: Several live routes are still searching for the right preserved quantity or witness measure, so invariant discovery remains a first-class discovery move.
\item Report: \texttt{C:\textbackslash\{\}files\textbackslash\{\}projects\textbackslash\{\}abducer\textbackslash\{\}abducer\textbackslash\{\}artifacts\textbackslash\{\}p-vs-np\textbackslash\{\}discovery-branch-executions\textbackslash\{\}discexec.p-vs-np.018.md}
\end{itemize}

\subsection*{disclayer.p-vs-np.pattern.invariant-discovery}
\begin{itemize}
\item Mode: discovery-candidate
\item Target: Invariant Discovery Prompt
\item Routing priority: 6.868
\item Linked queue items: 10
\item Headline claim: There exists an admissible MCSP witness family, realized as a family of instances equipped with a monotone witness measure, such that whenever A barely-nontrivial lower bound for an explicit problem Q in a weak model holds for MCSP in nonuniform-circuits, the object organizes the source-side mathematical content while preserving the invariant MCSP complexity profile in nonuniform-circuits and yields major consequences such as NP not in P/poly.
\item Rationale: Several live routes are still searching for the right preserved quantity or witness measure, so invariant discovery remains a first-class discovery move.
\item Report: \texttt{C:\textbackslash\{\}files\textbackslash\{\}projects\textbackslash\{\}abducer\textbackslash\{\}abducer\textbackslash\{\}artifacts\textbackslash\{\}p-vs-np\textbackslash\{\}discovery-branch-executions\textbackslash\{\}discexec.p-vs-np.018.md}
\end{itemize}

\subsection*{disclayer.p-vs-np.pattern.regime-localization}
\begin{itemize}
\item Mode: discovery-candidate
\item Target: Regime Localization Prompt
\item Routing priority: 6.758
\item Linked queue items: 10
\item Headline claim: There exists an admissible 'MCSP witness family' such that 'A hardness magnification theorem has the form:"A barely-nontrivial lower bound for an explicit problem Q in a weak model implies a major lower bound/separation in a strong model."Typically, Q is a meta-problem (e.g. a Gap-MCSP variant) and"major"means consequences such as NP not in P/poly.' can be organized into a route preserving 'MCSP complexity profile in nonuniform-circuits' and thereby supports 'In refined regimes, lower bounds for Gap-MCSP slightly beyond known frontiers (e.g. N\textasciicircum{} 1+varepsilon -type circuit lower bounds for suitable promise gaps and input length N=2\textasciicircum{}n ) imply major consequences such as NP not in P/poly.'.
\item Rationale: The frontier repeatedly crosses model or threshold boundaries, so discovery should search for the smallest regime in which the route first becomes true.
\item Report: \texttt{C:\textbackslash\{\}files\textbackslash\{\}projects\textbackslash\{\}abducer\textbackslash\{\}abducer\textbackslash\{\}artifacts\textbackslash\{\}p-vs-np\textbackslash\{\}discovery-branch-executions\textbackslash\{\}discexec.p-vs-np.020.md}
\end{itemize}

\subsection*{disclayer.p-vs-np.pattern.regime-localization}
\begin{itemize}
\item Mode: discovery-candidate
\item Target: Regime Localization Prompt
\item Routing priority: 6.758
\item Linked queue items: 10
\item Headline claim: There exists an admissible MCSP witness family, realized as a family of instances equipped with a monotone witness measure, such that whenever A barely-nontrivial lower bound for an explicit problem Q in a weak model holds for MCSP in nonuniform-circuits, the object organizes the source-side mathematical content while preserving the invariant MCSP complexity profile in nonuniform-circuits and yields major consequences such as NP not in P/poly.
\item Rationale: The frontier repeatedly crosses model or threshold boundaries, so discovery should search for the smallest regime in which the route first becomes true.
\item Report: \texttt{C:\textbackslash\{\}files\textbackslash\{\}projects\textbackslash\{\}abducer\textbackslash\{\}abducer\textbackslash\{\}artifacts\textbackslash\{\}p-vs-np\textbackslash\{\}discovery-branch-executions\textbackslash\{\}discexec.p-vs-np.020.md}
\end{itemize}

\subsection*{disclayer.p-vs-np.pattern.hybridization-of-partial-bridges}
\begin{itemize}
\item Mode: discovery-candidate
\item Target: Hybrid Bridge Synthesis Prompt
\item Routing priority: 6.728
\item Linked queue items: 10
\item Headline claim: There exists an admissible 'MCSP witness family' such that 'A hardness magnification theorem has the form:"A barely-nontrivial lower bound for an explicit problem Q in a weak model implies a major lower bound/separation in a strong model."Typically, Q is a meta-problem (e.g. a Gap-MCSP variant) and"major"means consequences such as NP not in P/poly.' can be organized into a route preserving 'MCSP complexity profile in nonuniform-circuits' and thereby supports 'In refined regimes, lower bounds for Gap-MCSP slightly beyond known frontiers (e.g. N\textasciicircum{} 1+varepsilon -type circuit lower bounds for suitable promise gaps and input length N=2\textasciicircum{}n ) imply major consequences such as NP not in P/poly.'.
\item Rationale: Some blocked routes sit at the intersection of multiple breakthrough families, which is evidence that one bridge object may be insufficient and a composed route should be explored.
\item Report: \texttt{C:\textbackslash\{\}files\textbackslash\{\}projects\textbackslash\{\}abducer\textbackslash\{\}abducer\textbackslash\{\}artifacts\textbackslash\{\}p-vs-np\textbackslash\{\}discovery-branch-executions\textbackslash\{\}discexec.p-vs-np.021.md}
\end{itemize}

\subsection*{disclayer.p-vs-np.pattern.hybridization-of-partial-bridges}
\begin{itemize}
\item Mode: discovery-candidate
\item Target: Hybrid Bridge Synthesis Prompt
\item Routing priority: 6.728
\item Linked queue items: 10
\item Headline claim: There exists an admissible MCSP witness family, realized as a family of instances equipped with a monotone witness measure, such that whenever A barely-nontrivial lower bound for an explicit problem Q in a weak model holds for MCSP in nonuniform-circuits, the object organizes the source-side mathematical content while preserving the invariant MCSP complexity profile in nonuniform-circuits and yields major consequences such as NP not in P/poly.
\item Rationale: Some blocked routes sit at the intersection of multiple breakthrough families, which is evidence that one bridge object may be insufficient and a composed route should be explored.
\item Report: \texttt{C:\textbackslash\{\}files\textbackslash\{\}projects\textbackslash\{\}abducer\textbackslash\{\}abducer\textbackslash\{\}artifacts\textbackslash\{\}p-vs-np\textbackslash\{\}discovery-branch-executions\textbackslash\{\}discexec.p-vs-np.021.md}
\end{itemize}

\subsection*{disclayer.p-vs-np.pattern.counterexample-led-concept-formation}
\begin{itemize}
\item Mode: discovery-candidate
\item Target: Counterexample-Led Concept Formation Prompt
\item Routing priority: 6.658
\item Linked queue items: 10
\item Headline claim: There exists an admissible MCSP witness family, realized as a family of instances equipped with a monotone witness measure, such that whenever A barely-nontrivial lower bound for an explicit problem Q in a weak model holds for MCSP in nonuniform-circuits, the object organizes the source-side mathematical content while preserving the invariant MCSP complexity profile in nonuniform-circuits and yields major consequences such as NP not in P/poly.
\item Rationale: The frontier is heavily blocked, so near-counterexamples and barrier analogs should be used to induce better concepts rather than only to reject claims.
\item Report: \texttt{C:\textbackslash\{\}files\textbackslash\{\}projects\textbackslash\{\}abducer\textbackslash\{\}abducer\textbackslash\{\}artifacts\textbackslash\{\}p-vs-np\textbackslash\{\}discovery-branch-executions\textbackslash\{\}discexec.p-vs-np.007.md}
\end{itemize}

\subsection*{disclayer.p-vs-np.pattern.route-decomposition}
\begin{itemize}
\item Mode: discovery-candidate
\item Target: Route Decomposition Prompt
\item Routing priority: 6.578
\item Linked queue items: 10
\item Headline claim: There exists an admissible MCSP witness family, realized as a family of instances equipped with a monotone witness measure, such that whenever A barely-nontrivial lower bound for an explicit problem Q in a weak model holds for MCSP in nonuniform-circuits, the object organizes the source-side mathematical content while preserving the invariant MCSP complexity profile in nonuniform-circuits and yields major consequences such as NP not in P/poly.
\item Rationale: High-pressure obstruction classes are carrying multiple inferential roles at once, so route decomposition is needed before more local search.
\item Report: \texttt{C:\textbackslash\{\}files\textbackslash\{\}projects\textbackslash\{\}abducer\textbackslash\{\}abducer\textbackslash\{\}artifacts\textbackslash\{\}p-vs-np\textbackslash\{\}discovery-branch-executions\textbackslash\{\}discexec.p-vs-np.009.md}
\end{itemize}

\subsection*{disclayer.p-vs-np.pattern.problem-representation-shift}
\begin{itemize}
\item Mode: discovery-candidate
\item Target: Problem Representation Shift Prompt
\item Routing priority: 6.548
\item Linked queue items: 10
\item Headline claim: There exists an admissible MCSP witness family, realized as a family of instances equipped with a monotone witness measure, such that whenever A barely-nontrivial lower bound for an explicit problem Q in a weak model holds for MCSP in nonuniform-circuits, the object organizes the source-side mathematical content while preserving the invariant MCSP complexity profile in nonuniform-circuits and yields major consequences such as NP not in P/poly.
\item Rationale: Several blocked routes look misrepresented rather than simply false, so representation change is itself a discovery pattern.
\item Report: \texttt{C:\textbackslash\{\}files\textbackslash\{\}projects\textbackslash\{\}abducer\textbackslash\{\}abducer\textbackslash\{\}artifacts\textbackslash\{\}p-vs-np\textbackslash\{\}discovery-branch-executions\textbackslash\{\}discexec.p-vs-np.011.md}
\end{itemize}

\subsection*{disclayer.p-vs-np.family.realization}
\begin{itemize}
\item Mode: discovery-candidate
\item Target: Carrier Realization Prompt
\item Routing priority: 6.443
\item Linked queue items: 9
\item Headline claim: If source theorem S holds in regime R, then one can construct an explicit carrier C with admissibility conditions A such that C preserves invariant I and supports a route-specific forcing lemma to target T.
\item Rationale: Turn a typed but schematic carrier into an explicit mathematical object family with admissibility conditions and a forcing theorem.
\item Report: \texttt{C:\textbackslash\{\}files\textbackslash\{\}projects\textbackslash\{\}abducer\textbackslash\{\}abducer\textbackslash\{\}artifacts\textbackslash\{\}p-vs-np\textbackslash\{\}discovery-branch-executions\textbackslash\{\}discexec.p-vs-np.004.md}
\end{itemize}

\subsection*{disclayer.p-vs-np.family.abstraction-search}
\begin{itemize}
\item Mode: discovery-candidate
\item Target: Abstraction Search Prompt
\item Routing priority: 6.443
\item Linked queue items: 9
\item Headline claim: There exists an admissible 'MCSP witness family' such that 'A hardness magnification theorem has the form:"A barely-nontrivial lower bound for an explicit problem Q in a weak model implies a major lower bound/separation in a strong model."Typically, Q is a meta-problem (e.g. a Gap-MCSP variant) and"major"means consequences such as NP not in P/poly.' can be organized into a route preserving 'MCSP complexity profile in nonuniform-circuits' and thereby supports 'In refined regimes, lower bounds for Gap-MCSP slightly beyond known frontiers (e.g. N\textasciicircum{} 1+varepsilon -type circuit lower bounds for suitable promise gaps and input length N=2\textasciicircum{}n ) imply major consequences such as NP not in P/poly.'.
\item Rationale: Find a sharper abstraction that turns the blocked route into a typed, inspectable theorem target.
\item Report: \texttt{C:\textbackslash\{\}files\textbackslash\{\}projects\textbackslash\{\}abducer\textbackslash\{\}abducer\textbackslash\{\}artifacts\textbackslash\{\}p-vs-np\textbackslash\{\}discovery-branch-executions\textbackslash\{\}discexec.p-vs-np.019.md}
\end{itemize}

\subsection*{disclayer.p-vs-np.family.abstraction-search}
\begin{itemize}
\item Mode: discovery-candidate
\item Target: Abstraction Search Prompt
\item Routing priority: 6.443
\item Linked queue items: 9
\item Headline claim: State a route-specific theorem that would resolve the blocked interface and make the next proof obligation materially sharper.
\item Rationale: Find a sharper abstraction that turns the blocked route into a typed, inspectable theorem target.
\item Report: \texttt{C:\textbackslash\{\}files\textbackslash\{\}projects\textbackslash\{\}abducer\textbackslash\{\}abducer\textbackslash\{\}artifacts\textbackslash\{\}p-vs-np\textbackslash\{\}discovery-branch-executions\textbackslash\{\}discexec.p-vs-np.019.md}
\end{itemize}

\subsection*{disclayer.p-vs-np.family.bridge-specificity}
\begin{itemize}
\item Mode: discovery-candidate
\item Target: Activation Theorem Prompt
\item Routing priority: 5.669
\item Linked queue items: 9
\item Headline claim: Target theorem T is fixed. The real theorem under search is the local activation theorem A that verifies T's hypotheses from current artifacts.
\item Rationale: Stop treating the target theorem as the open problem and relocate search to the bridge that would activate it.
\item Report: \texttt{C:\textbackslash\{\}files\textbackslash\{\}projects\textbackslash\{\}abducer\textbackslash\{\}abducer\textbackslash\{\}artifacts\textbackslash\{\}p-vs-np\textbackslash\{\}discovery-branch-executions\textbackslash\{\}discexec.p-vs-np.006.md}
\end{itemize}

\subsection*{disclayer.p-vs-np.family.transport}
\begin{itemize}
\item Mode: discovery-candidate
\item Target: Transport Principle Prompt
\item Routing priority: 5.610
\item Linked queue items: 8
\item Headline claim: For every admissible carrier C in the nonuniform-circuit regime, there is a transformed carrier C' in the formula-size regime preserving invariant I and retaining enough quantitative strength to imply target T.
\item Rationale: Find a theorem that preserves the relevant invariant across a cross-model move into formula-size.
\item Report: \texttt{C:\textbackslash\{\}files\textbackslash\{\}projects\textbackslash\{\}abducer\textbackslash\{\}abducer\textbackslash\{\}artifacts\textbackslash\{\}p-vs-np\textbackslash\{\}discovery-branch-executions\textbackslash\{\}discexec.p-vs-np.008.md}
\end{itemize}

\subsection*{disclayer.p-vs-np.family.barrier-aware-transfer}
\begin{itemize}
\item Mode: prove-candidate
\item Target: Barrier-Aware Transfer Prompt
\item Routing priority: 2.309
\item Linked queue items: 3
\item Headline claim: There exists an admissible 'unspecified witness family' such that 'There exist many problems L in combinatorial optimisation/graph theory such that L is NP-complete under polynomial-time many-one reductions.' can be organized into a route preserving 'consequence schema NP-complete-classification' and thereby supports 'Under"local"polynomial-time reductions (where each output bit is computable with very small time), MCSP is not NP-complete. More broadly, the paper develops evidence and conditional consequences suggesting NP-hardness of MCSP under various reductions would entail unexpected class separations.'.
\item Rationale: Design a transfer mechanism that survives known local-reduction barriers, or prove that no such mechanism exists under explicit assumptions.
\item Report: \texttt{C:\textbackslash\{\}files\textbackslash\{\}projects\textbackslash\{\}abducer\textbackslash\{\}abducer\textbackslash\{\}artifacts\textbackslash\{\}p-vs-np\textbackslash\{\}discovery-branch-executions\textbackslash\{\}discexec.p-vs-np.022.md}
\end{itemize}

\subsection*{disclayer.p-vs-np.family.barrier-aware-transfer}
\begin{itemize}
\item Mode: discovery-candidate
\item Target: Barrier-Aware Transfer Prompt
\item Routing priority: 2.309
\item Linked queue items: 3
\item Headline claim: Either construct a local reduction family G satisfying locality bounds L and preserving witness structure W, or prove that any such G would contradict barrier theorem B.
\item Rationale: Design a transfer mechanism that survives known local-reduction barriers, or prove that no such mechanism exists under explicit assumptions.
\item Report: \texttt{C:\textbackslash\{\}files\textbackslash\{\}projects\textbackslash\{\}abducer\textbackslash\{\}abducer\textbackslash\{\}artifacts\textbackslash\{\}p-vs-np\textbackslash\{\}discovery-branch-executions\textbackslash\{\}discexec.p-vs-np.022.md}
\end{itemize}

\subsection*{disclayer.p-vs-np.family.alignment}
\begin{itemize}
\item Mode: discovery-candidate
\item Target: Alignment Theorem Prompt
\item Routing priority: 2.222
\item Linked queue items: 3
\item Headline claim: Every admissible carrier C from the local route induces an object D in the target theorem's class, and D preserves the quantitative property needed to activate the target theorem.
\item Rationale: Identify or prove an equivalence between the current carrier class and the object class required by the target theorem.
\item Report: \texttt{C:\textbackslash\{\}files\textbackslash\{\}projects\textbackslash\{\}abducer\textbackslash\{\}abducer\textbackslash\{\}artifacts\textbackslash\{\}p-vs-np\textbackslash\{\}discovery-branch-executions\textbackslash\{\}discexec.p-vs-np.012.md}
\end{itemize}

\subsection*{disclayer.p-vs-np.family.explicitness}
\begin{itemize}
\item Mode: prove-candidate
\item Target: Explicitness Upgrade Prompt
\item Routing priority: 2.107
\item Linked queue items: 3
\item Headline claim: There exists an admissible 'PIT certificate family' such that 'A set Hsubseteq mathbb F \textasciicircum{}n is a hitting set for a class mathcal C of n -variate circuits if for every nonzero polynomial p computed by some circuit in mathcal C , there exists mathbf a in H with p(mathbf a )neq 0.' can be organized into a route preserving 'PIT semantic object' and thereby supports 'If mathrm PIT in Pclass (or even suitable infinitely-often derandomisations in nondeterministic subexponential time, in standard variants), then at least one of the following holds: NEXP not in P/poly quad or quad Permanent does not have polynomial-size arithmetic circuits.'.
\item Rationale: Upgrade a semantic or existential algebraic object into an explicit algorithmic derandomization premise.
\item Report: \texttt{C:\textbackslash\{\}files\textbackslash\{\}projects\textbackslash\{\}abducer\textbackslash\{\}abducer\textbackslash\{\}artifacts\textbackslash\{\}p-vs-np\textbackslash\{\}discovery-branch-executions\textbackslash\{\}discexec.p-vs-np.023.md}
\end{itemize}

\subsection*{disclayer.p-vs-np.family.explicitness}
\begin{itemize}
\item Mode: discovery-candidate
\item Target: Explicitness Upgrade Prompt
\item Routing priority: 2.107
\item Linked queue items: 3
\item Headline claim: If semantic object C exists for the relevant circuit class, then one can uniformly and efficiently extract an explicit structure G that yields deterministic PIT for that class.
\item Rationale: Upgrade a semantic or existential algebraic object into an explicit algorithmic derandomization premise.
\item Report: \texttt{C:\textbackslash\{\}files\textbackslash\{\}projects\textbackslash\{\}abducer\textbackslash\{\}abducer\textbackslash\{\}artifacts\textbackslash\{\}p-vs-np\textbackslash\{\}discovery-branch-executions\textbackslash\{\}discexec.p-vs-np.023.md}
\end{itemize}

\subsection*{disclayer.p-vs-np.family.constructive-instantiation}
\begin{itemize}
\item Mode: discovery-candidate
\item Target: Constructive Instantiation Prompt
\item Routing priority: 1.966
\item Linked queue items: 3
\item Headline claim: From abstract schema S, construct an explicit family of natural instances E\_n satisfying properties P and realizing target theorem T.
\item Rationale: Turn an abstract schema into explicit natural example families or prove that the schema cannot instantiate naturally in the current way.
\item Report: \texttt{C:\textbackslash\{\}files\textbackslash\{\}projects\textbackslash\{\}abducer\textbackslash\{\}abducer\textbackslash\{\}artifacts\textbackslash\{\}p-vs-np\textbackslash\{\}discovery-branch-executions\textbackslash\{\}discexec.p-vs-np.014.md}
\end{itemize}

\subsection*{disclayer.p-vs-np.family.semantic-bridge}
\begin{itemize}
\item Mode: discovery-candidate
\item Target: Semantic Bridge Prompt
\item Routing priority: 1.946
\item Linked queue items: 3
\item Headline claim: Local definitional object D has property P if and only if global characterization condition Q holds, with both directions proved at the right level of abstraction.
\item Rationale: Bridge a local definitional object to a global characterization theorem.
\item Report: \texttt{C:\textbackslash\{\}files\textbackslash\{\}projects\textbackslash\{\}abducer\textbackslash\{\}abducer\textbackslash\{\}artifacts\textbackslash\{\}p-vs-np\textbackslash\{\}discovery-branch-executions\textbackslash\{\}discexec.p-vs-np.015.md}
\end{itemize}

\subsection*{disclayer.p-vs-np.family.categorical-reframing}
\begin{itemize}
\item Mode: discovery-candidate
\item Target: Categorical Reframing Prompt
\item Routing priority: 0.910
\item Linked queue items: 12
\item Headline claim: There exists a mediating object built from 'MCSP witness family' such that the route from 'A hardness magnification theorem has the form:"A barely-nontrivial lower bound for an explicit problem Q in a weak model implies a major lower bound/separation in a strong model."Typically, Q is a meta-problem (e.g. a Gap-MCSP variant) and"major"means consequences such as NP not in P/poly.' to 'In refined regimes, lower bounds for Gap-MCSP slightly beyond known frontiers (e.g. N\textasciicircum{} 1+varepsilon -type circuit lower bounds for suitable promise gaps and input length N=2\textasciicircum{}n ) imply major consequences such as NP not in P/poly.' factors through typed morphisms preserving 'MCSP complexity profile in nonuniform-circuits' and making the relevant square commute.
\item Rationale: Recast the blocked frontier as a diagrammatic composition problem: identify objects, morphisms, mediating objects, and the exact commutativity or factorization obligation.
\item Report: \texttt{C:\textbackslash\{\}files\textbackslash\{\}projects\textbackslash\{\}abducer\textbackslash\{\}abducer\textbackslash\{\}artifacts\textbackslash\{\}p-vs-np\textbackslash\{\}discovery-branch-executions\textbackslash\{\}discexec.p-vs-np.024.md}
\end{itemize}

\subsection*{disclayer.p-vs-np.family.categorical-reframing}
\begin{itemize}
\item Mode: discovery-candidate
\item Target: Categorical Reframing Prompt
\item Routing priority: 0.910
\item Linked queue items: 12
\item Headline claim: There exists a mediating object built from 'MCSP witness family' such that the route from 'A hardness magnification theorem has the form:"A barely-nontrivial lower bound for an explicit problem Q in a weak model implies a major lower bound/separation in a strong model."Typically, Q is a meta-problem (e.g. a Gap-MCSP variant) and"major"means consequences such as NP not in P/poly.' to 'In refined regimes, lower bounds for Gap-MCSP slightly beyond known frontiers (e.g. N\textasciicircum{} 1+varepsilon -type circuit lower bounds for suitable promise gaps and input length N=2\textasciicircum{}n ) imply major consequences such as NP not in P/poly.' factors through typed morphisms preserving 'MCSP complexity profile in nonuniform-circuits' and making the relevant square commute.
\item Rationale: Recast the blocked frontier as a diagrammatic composition problem: identify objects, morphisms, mediating objects, and the exact commutativity or factorization obligation.
\item Report: \texttt{C:\textbackslash\{\}files\textbackslash\{\}projects\textbackslash\{\}abducer\textbackslash\{\}abducer\textbackslash\{\}artifacts\textbackslash\{\}p-vs-np\textbackslash\{\}discovery-branch-executions\textbackslash\{\}discexec.p-vs-np.024.md}
\end{itemize}

\subsection*{disclayer.p-vs-np.family.equivalence-coherence-reframing}
\begin{itemize}
\item Mode: discovery-candidate
\item Target: Equivalence / Coherence Reframing Prompt
\item Routing priority: 0.880
\item Linked queue items: 12
\item Headline claim: There exists an equivalence witness identifying 'MCSP witness family' with the target-side object class relevant to 'In refined regimes, lower bounds for Gap-MCSP slightly beyond known frontiers (e.g. N\textasciicircum{} 1+varepsilon -type circuit lower bounds for suitable promise gaps and input length N=2\textasciicircum{}n ) imply major consequences such as NP not in P/poly.', and 'MCSP complexity profile in nonuniform-circuits' transports coherently along that witness.
\item Rationale: Recast the blocked frontier as an identity, equivalence, transport, or coherence problem rather than as a missing brute-force bridge.
\item Report: \texttt{C:\textbackslash\{\}files\textbackslash\{\}projects\textbackslash\{\}abducer\textbackslash\{\}abducer\textbackslash\{\}artifacts\textbackslash\{\}p-vs-np\textbackslash\{\}discovery-branch-executions\textbackslash\{\}discexec.p-vs-np.025.md}
\end{itemize}

\subsection*{disclayer.p-vs-np.family.equivalence-coherence-reframing}
\begin{itemize}
\item Mode: discovery-candidate
\item Target: Equivalence / Coherence Reframing Prompt
\item Routing priority: 0.880
\item Linked queue items: 12
\item Headline claim: There exists an equivalence witness identifying 'MCSP witness family' with the target-side object class relevant to 'In refined regimes, lower bounds for Gap-MCSP slightly beyond known frontiers (e.g. N\textasciicircum{} 1+varepsilon -type circuit lower bounds for suitable promise gaps and input length N=2\textasciicircum{}n ) imply major consequences such as NP not in P/poly.', and 'MCSP complexity profile in nonuniform-circuits' transports coherently along that witness.
\item Rationale: Recast the blocked frontier as an identity, equivalence, transport, or coherence problem rather than as a missing brute-force bridge.
\item Report: \texttt{C:\textbackslash\{\}files\textbackslash\{\}projects\textbackslash\{\}abducer\textbackslash\{\}abducer\textbackslash\{\}artifacts\textbackslash\{\}p-vs-np\textbackslash\{\}discovery-branch-executions\textbackslash\{\}discexec.p-vs-np.025.md}
\end{itemize}

\end{document}
